% !TeX encoding = UTF-8
% !TeX program = pdflatex
% ---------------------------------------------------------------------------
% IEEE TETCI (Transactions on Emerging Topics in Computational Intelligence)
%
% Usage:
%   1. Requires IEEEtran.cls and IEEEtran.bst (copy from templates/ieee-latex/)
%   2. Compile: pdflatex → bibtex → pdflatex → pdflatex
%
% Journal: IEEE Transactions on Emerging Topics in Computational Intelligence
% Publisher: IEEE Computational Intelligence Society (CIS)
% ISSN: 2471-285X
% First published: 2017
% Scope: Neural networks, evolutionary computation, fuzzy systems, and
%        emerging computational intelligence topics
% ---------------------------------------------------------------------------

\documentclass[journal,onecolumn,draftcls]{IEEEtran}
\IEEEoverridecommandlockouts

% --- Final version ---
% \documentclass[journal]{IEEEtran}
% \IEEEoverridecommandlockouts

\usepackage{cite}
\usepackage{amsmath,amssymb,amsfonts}
\usepackage{algorithmic}
\usepackage{algorithm}
\usepackage{graphicx}
\usepackage{textcomp}
\usepackage{xcolor}
\usepackage{booktabs}
\usepackage{multirow}
\usepackage{hyperref}
\usepackage{subcaption}

\hypersetup{
  colorlinks=true,
  linkcolor=blue,
  citecolor=blue,
  urlcolor=blue,
}

\def\BibTeX{{\rm B\kern-.05em{\sc i\kern-.025em b}\kern-.08em
    T\kern-.1667em\lower.7ex\hbox{E}\kern-.125emX}}

% --- TETCI Journal Metadata ---
% \markboth{IEEE Transactions on Emerging Topics in Computational Intelligence, Vol. XX, No. X, Month 2026}%
% {Author \MakeLowercase{\textit{et al.}}: Short Title Here}

\begin{document}

% ===========================================================================
% TITLE
% ===========================================================================
\title{Your Paper Title: [Method Name] for [Computational Intelligence Task]}

\author{
  \IEEEauthorblockN{Author One\IEEEauthorrefmark{1},
                    Author Two\IEEEauthorrefmark{2},
                    Author Three\IEEEauthorrefmark{1}}
  \IEEEauthorblockA{
    \IEEEauthorrefmark{1}Department of XXX, University of YYY, City, Country \\
    \IEEEauthorrefmark{2}Institute of ZZZ, Academy of Sciences, City, Country \\
    E-mail: \{email1, email2\}@university.edu
  }
}

\maketitle

% ===========================================================================
% ABSTRACT
% ===========================================================================
\begin{abstract}

% TETCI abstracts: 150-250 words, single paragraph.
% Must clearly convey the computational intelligence contribution.
%
% TETCI emphasizes: (1) emerging/novel CI methods,
% (2) rigorous theoretical or empirical validation,
% (3) clear motivation for WHY a new CI method is needed.

Your abstract here. Structure:
\begin{itemize}
  \item \textbf{Context:} What CI problem does this paper address?
  \item \textbf{Gap:} Why existing CI methods are insufficient.
  \item \textbf{Method:} Your CI approach -- First... Second... Finally...
  \item \textbf{Results:} Quantitative evidence of effectiveness.
  \item \textbf{Implication:} What this means for the CI community.
\end{itemize}

\end{abstract}

% ===========================================================================
% KEYWORDS
% ===========================================================================
\begin{IEEEkeywords}
  Computational intelligence, [your CI paradigm, e.g., evolutionary computation,
  neural networks, fuzzy systems], [your technique, e.g., multi-objective
  optimization, deep reinforcement learning, transfer learning],
  [application domain, e.g., industrial optimization, medical diagnosis].
\end{IEEEkeywords}

% ===========================================================================
% 1. INTRODUCTION
% ===========================================================================
\section{Introduction}
\label{sec:introduction}

% TETCI Introduction (5 paragraphs, ~2 pages journal style):

% Para 1: CI context and practical motivation.
%   Why this problem matters from a computational intelligence perspective.
%   Connect to real-world applications that require CI solutions.

% Para 2: Current CI approaches.
%   Review 2-3 major CI paradigms applied to this problem:
%   - Evolutionary computation / swarm intelligence
%   - Neural networks / deep learning
%   - Fuzzy systems / rough sets
%   - Hybrid CI approaches
%   Discuss what each paradigm contributes and where it falls short.

% Para 3: Remaining gap.
%   Must contain a "because" clause specific to the CI challenge.
%   Examples of CI-specific gaps:
%   - "Existing evolutionary algorithms converge slowly because the search space
%      grows exponentially with [problem characteristic]."
%   - "Current neural network approaches lack interpretability because [reason],
%      which is critical for [application]."
%   - "Fuzzy rule-based systems fail to handle [data characteristic] because
%      their membership functions cannot adapt to [condition]."

% Para 4: Proposed method teaser.
%   3-4 sentences explaining the CI innovation and core intuition.

% Para 5: Contributions (4-5 bullets for TETCI journal).

The main contributions of this paper are:
\begin{enumerate}
  \item We identify [specific limitation] in existing [CI paradigm] for
    [problem class], and we provide quantitative evidence of this limitation
    (Section~\ref{sec:motivation}).
  \item We propose [Method Name], which [core CI mechanism] to address
    [specific CI challenge] (Section~\ref{sec:method}).
  \item We provide [theoretical analysis / convergence proof / complexity
    analysis] demonstrating that [property] holds for the proposed method
    (Section~\ref{sec:theory}).
  \item Through experiments on [benchmark suites / real-world problems],
    [Method] achieves [key result], outperforming [N] state-of-the-art
    CI methods by [X]\% on [metric] (Section~\ref{sec:experiments}).
  \item We release [code / benchmark / dataset] to facilitate reproducible
    CI research.
\end{enumerate}

% ===========================================================================
% 2. RELATED WORK
% ===========================================================================
\section{Related Work}
\label{sec:related}

% Organize by CI paradigm, not paper-by-paper.
% TETCI expects breadth across CI subfields.

\subsection{[CI Paradigm A]}
\label{sec:related-paradigmA}

% e.g., Evolutionary Computation approaches to this problem

\subsection{[CI Paradigm B]}
\label{sec:related-paradigmB}

% e.g., Neural Network / Deep Learning approaches

\subsection{[CI Paradigm C]}
\label{sec:related-paradigmC}

% e.g., Fuzzy Logic / Rough Set approaches (if applicable)

\subsection{Hybrid and Emerging CI Approaches}
\label{sec:related-hybrid}

% Discuss hybrid methods (e.g., neuro-evolution, fuzzy-neural, etc.)
% This is especially relevant for an "Emerging Topics" journal.

\subsection{Positioning of This Work}
\label{sec:related-positioning}

% Explicitly differentiate: what CI paradigm(s) do you combine or advance?
% How is your work different from ALL reviewed approaches?

% ===========================================================================
% 3. PRELIMINARIES AND PROBLEM FORMULATION
% ===========================================================================
\section{Preliminaries and Problem Formulation}
\label{sec:preliminaries}

% TETCI papers often include formal definitions.
% For CI papers, define:
% - The optimization/search problem
% - The decision space and objective space
% - Constraints and assumptions
% - Performance criteria

\subsection{Background}
\label{sec:prelim-bg}

% Define relevant CI concepts needed to understand your method.
% Example: if proposing a new evolutionary operator, explain the base EA first.

\subsection{Problem Definition}
\label{sec:prelim-problem}

% Formal mathematical definition of the problem.
% For optimization: $\min_{\mathbf{x} \in \Omega} f(\mathbf{x})$
% For learning: define input space $\mathcal{X}$, output space $\mathcal{Y}$,
%   loss function, hypothesis class.

% ===========================================================================
% 4. METHOD
% ===========================================================================
\section{Method}
\label{sec:method}

% Opening paragraph (3-5 sentences): overall framework → what each of the
% three innovations solves → how they collaborate. No "Overview" subsection.
% Explain the CI design philosophy: WHY this approach, not just WHAT.

\begin{figure}[t]
  \centering
  \fbox{\parbox{0.85\linewidth}{\centering
  \vspace{3cm}
  {\large [Fig. 1: Overview of the proposed CI method.]}\\
  \vspace{0.3cm}
  {\footnotesize The framework consists of: (a) [Initialization / encoding],
   (b) [Core CI mechanism],
   (c) [Adaptation / learning strategy],
   (d) [Output / decision module].}
  \vspace{3cm}}}
  \caption{Overview of the proposed [Method Name].}
  \label{fig:architecture}
\end{figure}

\subsection{<Innovation 1 Name>}
\label{sec:method-innov1}

% Motivation → Design → Formula/Mechanism → Effect
% Include equations and/or pseudocode as needed.

\begin{algorithm}[t]
  \caption{[Core Component] of the Proposed Method}
  \label{alg:core}
  \begin{algorithmic}[1]
    \STATE \textbf{Input:} [problem parameters], [hyperparameters]
    \STATE Initialize [population / network / rule base]
    \FOR{each [iteration / generation / epoch]}
      \STATE [Core operation step 1]
      \STATE [Core operation step 2]
      \STATE [Adaptation / update step]
    \ENDFOR
    \STATE \textbf{Output:} [solution / model / decision]
  \end{algorithmic}
\end{algorithm}

\subsection{<Innovation 2 Name>}
\label{sec:method-innov2}

% Motivation → Design → Formula/Mechanism → Effect

\subsection{<Innovation 3 Name>}
\label{sec:method-innov3}

% Motivation → Design → Formula/Mechanism → Effect

\subsection{Implementation Details}
\label{sec:method-implementation}

% Framework, optimizer, LR schedule, complexity analysis, etc.
% TETCI strongly values computational complexity analysis.

% ===========================================================================
% 5. THEORETICAL ANALYSIS (if applicable)
% ===========================================================================
% \section{Theoretical Analysis}
% \label{sec:theory}

% ===========================================================================
% 6. EXPERIMENTS
% ===========================================================================
\section{Experiments}
\label{sec:experiments}

\subsection{Datasets and Metrics}
\label{sec:exp-data}

% Benchmark problems, datasets, metrics with direction.
% TETCI requires rigorous statistics: Wilcoxon, Friedman test, effect size.

\begin{table}[t]
  \centering
  \caption{Summary of benchmark problems / datasets.}
  \label{tab:benchmarks}
  \begin{tabular}{lccc}
    \toprule
    \textbf{Problem} & \textbf{Dim} & \textbf{Characteristics} &
    \textbf{Source} \\
    \midrule
    F1--F10 & 30/50/100 & Unimodal & CEC 2024 \cite{ref1} \\
    F11--F20 & 30/50/100 & Multimodal & CEC 2024 \cite{ref1} \\
    Real-world A & variable & [description] & [source] \\
    \bottomrule
  \end{tabular}
\end{table}

\subsection{Experimental Setup}
\label{sec:exp-setup}

% All compared methods with citations and parameter settings.
% Programming language, libraries, hardware, number of runs, stopping criteria.
% TETCI requires rigorous statistics: Wilcoxon, Friedman, critical difference diagrams.

\subsection{Comparison with State-of-the-Art and Classical Methods}
\label{sec:exp-main}

% Results tables with statistical annotations. \pm for std.
% Include classical CI baselines and recent SOTA.

\subsection{Ablation Study}
\label{sec:exp-ablation}

% Numbered list 1), 2), 3)... analyzing each component / design choice.
% May also include: parameter sensitivity, scalability analysis.

\subsection{Visualization}
\label{sec:exp-visual}

% Convergence curves, critical difference diagrams, Pareto fronts, etc.

% ===========================================================================
% 7. DISCUSSION
% ===========================================================================
\section{Discussion}
\label{sec:discussion}

\subsection{Why the Proposed CI Approach Works}
\label{sec:discuss-why}

\subsection{Limitations and When the Method May Fail}
\label{sec:discuss-limitations}

\subsection{Future Research Directions}
\label{sec:discuss-future}

% ===========================================================================
% 8. CONCLUSION
% ===========================================================================
\section{Conclusion}
\label{sec:conclusion}

% Paragraph 1: Full summary (problem → method → key results with numbers)
% Paragraph 2: Limitations + future work (specific limitations + actionable directions)

% ===========================================================================
% ACKNOWLEDGMENTS
% ===========================================================================
\section*{Acknowledgments}

% Funding sources, computational resources, helpful discussions.

% ===========================================================================
% APPENDICES (optional)
% ===========================================================================
% \appendices
% \section{Detailed Derivation of [Theorem/Equation]}
% \section{Additional Experimental Results}
% \section{Parameter Settings for Compared Methods}

% ===========================================================================
% REFERENCES
% ===========================================================================
{
  % TETCI uses numeric IEEE citation: [1], [2], ...
  % Include ~40-70 references spanning:
  % - Classic CI papers (Goldberg, Holland, Kennedy & Eberhart, Zadeh, etc.)
  % - Recent CI advances (deep neuroevolution, quality-diversity, etc.)
  % - Application-domain papers
  % - Benchmark/competition papers (CEC, GECCO proceedings)

  \bibliographystyle{IEEEtran}
  \bibliography{references}
}

\end{document}
